Compactness has been a legal criterion for redistricting in the United States since
the nineteenth century, yet it remains operationally undefined in federal law.
The Supreme Court held in \textit{Shaw v.\ Reno}, 509 U.S.\ 630 (1993), that
districts with shapes ``so bizarre'' as to be inexplicable except as racial
classifications trigger strict scrutiny under the Equal Protection Clause.
Justice O'Connor's majority opinion described districts ``so dramatically irregular
that, on its face, it rationally cannot be understood as anything other than an
effort to segregate voters by race''---language that implicitly invokes compactness
without defining it mathematically~\citep{shaw1993}.
\textit{Rucho v.\ Common Cause}, 588 U.S.\ 684 (2019), confirmed that partisan
gerrymandering claims are non-justiciable under federal law but explicitly noted
that state-law compactness standards remain enforceable~\citep{rucho2019}.
State constitutions in at least 22 states mandate compactness as a redistricting
criterion~\citep{ncsl2021criteria}.

Despite this legal backdrop, the redistricting literature has produced at least
six distinct mathematical operationalisations of compactness, each capturing
a different geometric property:
\begin{itemize}
  \item \textbf{Polsby-Popper (PP)}: area-to-perimeter ratio, penalising jagged
    boundaries~\citep{polsby1991third}.
  \item \textbf{Reock}: ratio of district area to its minimum bounding circle,
    measuring how round the district is in the broadest sense~\citep{reock1961note}.
  \item \textbf{Convex Hull Ratio (CH)}: ratio of district area to convex hull area,
    detecting non-convex ``tentacle'' geometries~\citep{young1988measuring}.
  \item \textbf{Schwartzberg (S)}: perimeter-to-circle-equivalent ratio,
    preferred in some state statutes~\citep{schwartzberg1966reapportionment}.
  \item \textbf{Length-Width Ratio (LW)}: minimum bounding rectangle elongation,
    used to identify linear districts~\citep{niemi1990measuring}.
  \item \textbf{Population-Weighted Compactness (PWC)}: moment-of-inertia metric
    weighting by tract population, measuring representational compactness~\citep{fryer2011measuring}.
\end{itemize}

No prior work has systematically compared all six metrics across a common algorithmic
redistricting platform with fixed census data and reproducible seeds. The \textsc{bisect}
redistricting system~\citep{bisect2024} provides that platform: four distinct
structure algorithms (standard-bisect, prime-factor/ApportionRegions, ratio-optimal/GeoSection,
moving-knife/MKA) each generate a complete set of districts for NC, WI, and TX,
enabling controlled comparison across 24 algorithm-state pairs.

\subsection{Contribution}

This paper makes three contributions:
\begin{enumerate}
  \item \textbf{Taxonomy}: We organise the six metrics into two families---area-based
    (Reock, CH) and perimeter-based (PP, S)---plus one representational metric (PWC),
    and characterise the within-family and across-family correlation structure.
  \item \textbf{Correlation analysis}: We report the full $6\times6$ Pearson
    correlation matrix across all districts produced by all four algorithms for
    NC, WI, and TX, providing the first empirical estimate of metric redundancy
    under controlled algorithmic redistricting.
  \item \textbf{Decision table}: We provide a practitioner guide mapping each
    metric to its primary legal use case, recommended court contexts, and
    \textsc{bisect} CLI flag.
\end{enumerate}

This paper is the foundational reference for K.1--K.7, each of which studies
one metric in depth.

\subsection{Data and Reproducibility}

All results use 2020 Census TIGER/Line tract geometries and PL 94-171
redistricting-file population counts.
Single-seed results (seed\,$=42$) are marked with~$^\dagger$.
Implementation is in \texttt{crates/bisect-analysis/src/compactness.rs}.
All six metrics are produced by:
\begin{verbatim}
bisect label-analyze <label> --year 2020 --types all
\end{verbatim}
